\section{A complete exponential expansion for the angular zero}
\label{angular:sec:main}

For integers $a\geq1$ and real $b>0$, write
\[
 H_j^{(b)}=\sum_{m=1}^j m^{-b},\qquad
 F_{a,b}(z)=\Li_{a,b}(z,1)
   =\sum_{n=2}^{\infty}\frac{H_{n-1}^{(b)}}{n^a}z^n.
\]
The signed-density theorem in the pinned manuscript
\cite[\texttt{signed:thm:unique-zero}]{ProveItSnapshot}
proves that $\Imag F_{a,b}(e^{i\theta})$ has exactly one zero
$\theta_{a,b}$ in the open upper semicircle, that it is simple, and
that $0<\theta_{a,b}<\pi/2$.  Its proof uses a zero-mass density
with one sign change and a strictly ordered Cauchy kernel.  The source
also proves a four-term expansion as $a\to\infty$, uniformly for
$b\geq1$, and explicitly leaves a systematic further expansion open.

We give every coefficient, with a remainder at any prescribed
exponential scale.  The result is uniform throughout $b>0$.  Only
the elementary bound
\begin{equation}\label{angular:eq:hbound}
 1\leq H_j^{(b)}\leq j\qquad(j\geq1,\ b>0)
\end{equation}
is needed for this stronger uniformity.  The expansion is complete in the asymptotic sense. Theorem
\ref{angular:thm:divergence} proves that its infinite untruncated sum
diverges at every fixed $a,b$, even after equal bases are combined.

\subsection{The locally finite scale and its coefficients}

Let $\mathscr M$ be the set of nonzero finitely supported sequences
$\mathbf m=(m_3,m_4,\ldots)$ of nonnegative integers. Put
\[
 |\mathbf m|=\sum_{n\geq3}m_n,
 \qquad q(\mathbf m)=\prod_{n\geq3}(2/n)^{m_n}.
\]
For every $0<\eta<1$, only finitely many $\mathbf m$ have
$q(\mathbf m)>\eta$. Indeed, each occupied index satisfies
$n<2/\eta$, and $(2/3)^{|\mathbf m|}\geq q(\mathbf m)>\eta$ bounds
the total multiplicity. Distinct multisets can give the same rational
base: those coefficients must be added, not interpreted as distinct
asymptotic scales.

Define the analytic germ
\begin{equation}\label{angular:eq:phi}
 \phi_n(u)=-\frac{u\sin(n\pi/2-nu)}{\sin(2u)}\qquad(n\geq3),
\end{equation}
with its removable value at $u=0$. Its Taylor coefficients are rational.
For $k=|\mathbf m|$ define
\begin{equation}\label{angular:eq:coefficient}
 C_{\mathbf m}(b)=
 \frac{(k-1)!}{\prod_{n\geq3}m_n!}
 [u^{k-1}]\prod_{n\geq3}\phi_n(u)^{m_n}
 \prod_{n\geq3}\bigl(H_{n-1}^{(b)}\bigr)^{m_n}.
\end{equation}
Thus $C_{\mathbf m}(b)$ is an explicitly given rational polynomial
in finitely many generalized harmonic numbers. In particular it is
rational when $b$ is a positive integer.

\begin{theorem}[All exponential orders, uniformly in the second index]
\label{angular:thm:allorders}
For every fixed $0<\eta<1$, as the integer $a\to\infty$,
\begin{equation}\label{angular:eq:allorders}
 \delta_{a,b}:=\frac\pi2-\theta_{a,b}
 =\sum_{\substack{\mathbf m\in\mathscr M\\q(\mathbf m)>\eta}}
 C_{\mathbf m}(b)q(\mathbf m)^a+O_\eta(\eta^a),
\end{equation}
uniformly for all real $b>0$. The constant and the threshold for $a$
depend on $\eta$, but not on $b$. The coefficient vanishes whenever
the total multiplicity $\sum_{n\text{ odd}}m_n$ is even.
\end{theorem}

\begin{proof}
For $a\geq4$, the Fourier series and its first angular derivative
converge absolutely, uniformly under \eqref{angular:eq:hbound}.
The equation for the zero, normalized by $2^a$, is
\begin{equation}\label{angular:eq:fourier}
 S_{a,b}(\delta)=\sin(2\delta)+
 \sum_{n\geq3}H_{n-1}^{(b)}(2/n)^a
      \sin(n\pi/2-n\delta)=0.
\end{equation}
The $n=3$ term at $\delta=0$ is
$-H_2^{(b)}(2/3)^a$; the remaining nonzero terms there have
size $O((2/5)^a)$ uniformly in $b>0$.  On
$0\leq\delta\leq2(2/3)^a$, differentiating the series gives
\[
 S_{a,b}'(\delta)=2+O((2/3)^a).
\]
For example, the differentiated tail is bounded by
\[
 \sum_{n\geq3}n(n-1)(2/n)^a
 \leq (2/3)^a\sum_{n\geq3}n^2(3/n)^4,
\]
whose last sum is finite. The same elementary estimate bounds the
undifferentiated tails. Since $1\leq H_2^{(b)}\leq2$,
$S_{a,b}(0)<0<S_{a,b}(2(2/3)^a)$ for all sufficiently large $a$.
Consequently the zero in this interval is unique and
\begin{equation}\label{angular:eq:smallroot}
 \delta_{a,b}=O((2/3)^a),\qquad
 S_{a,b}'(\delta)\geq1
\end{equation}
there, uniformly in $b$. The source's global uniqueness identifies
this local root with $\pi/2-\theta_{a,b}$.

If $\eta\geq2/3$, the sum in \eqref{angular:eq:allorders} is empty
and \eqref{angular:eq:smallroot} proves the assertion.  Otherwise
take $N=\lceil2/\eta\rceil-1\geq3$ and keep only the modes
$3\leq n\leq N$ in \eqref{angular:eq:fourier}. On the real root
interval the omitted tail is bounded by
\begin{align*}
 \sum_{n>N}(n-1)(2/n)^a
 &\leq (2/(N+1))^a
       \sum_{n>N}(n-1)((N+1)/n)^4\\
 &=O_\eta(\eta^a).
\end{align*}
The derivative bounds give a unique small root $\delta_N$ of this
finite equation and
\begin{equation}\label{angular:eq:finiteerror}
 \delta_{a,b}-\delta_N=O_\eta(\eta^a).
\end{equation}
One can apply the mean value theorem between the two roots; both
are in a common interval where the derivative of the finite
equation is bounded below by $1$.

Now regard $z_3,\ldots,z_N$ as independent complex variables. The
finite equation
\[
 \sin(2\delta)+\sum_{n=3}^Nz_n\sin(n\pi/2-n\delta)=0
\]
has a holomorphic solution $\delta=\Psi_N(\mathbf z)$ near
$\mathbf z=0$, by the implicit function theorem with derivative $2$.
Dividing by $\sin(2\delta)/\delta$ gives
\[
 \delta=\sum_{n=3}^Nz_n\phi_n(\delta).
\]
Apply the finite-variable Lagrange--B\"urmann formula, or insert an
auxiliary scalar multiplying all $z_n$ and use the ordinary formula.
The coefficient of $\prod z_n^{m_n}$ is precisely
\[
 \frac{(k-1)!}{\prod m_n!}
 [u^{k-1}]\prod\phi_n(u)^{m_n}.
\]
This is a finite-variable use of classical Lagrange inversion;
Dirichlet-series inversion has a wider literature, for example
Kuznetsov~\cite{Kuznetsov2019}, but no infinite-variable convergence
theorem is being assumed here.

Choose $K$ so that $(2/3)^{K+1}\leq\eta$. Truncate the convergent
Taylor expansion of $\Psi_N$ at total degree $K$, then substitute
$z_n=H_{n-1}^{(b)}(2/n)^a$. Its Taylor remainder is
$O_\eta((2/3)^{(K+1)a})=O_\eta(\eta^a)$, uniformly by the finite
bounds $H_{n-1}^{(b)}\leq N$. Every discarded monomial whose
base is at most $\eta$ also has size $O_\eta(\eta^a)$.
Every multiset with base greater than $\eta$ is supported on
$\{3,\ldots,N\}$ and has total degree at most $K$. Combining
this observation with \eqref{angular:eq:finiteerror} proves
\eqref{angular:eq:allorders}.

Finally, $\phi_n$ is even in $u$ for odd $n$ and odd in $u$ for
even $n$. Its product can have a coefficient at $u^{k-1}$ only if
$\sum_{n\text{ even}}m_n\equiv k-1\pmod2$, equivalently if
$\sum_{n\text{ odd}}m_n$ is odd. This proves the vanishing assertion.
\end{proof}

\begin{remark}[Why the proof truncates the Fourier modes first]
The original Fourier series does not converge absolutely in a complex
neighborhood of a real angular argument: high frequencies grow on
one side of that neighborhood. An unqualified infinite-dimensional
analytic implicit-function argument would therefore leave a real gap.
The finite cutoff followed by a real derivative estimate avoids that
gap and proves exactly the asymptotic assertion stated above. Theorem~\ref{angular:thm:divergence} proves that the formal sum over
all multisets, and even the sum grouped by equal bases, diverges.
\end{remark}

\subsection{Ten terms and the next signed remainder}

Write $h_j=H_j^{(b)}$. Let $T_{10}(a,b)$ be the sum of $c(q)q^a$
over the ten rows in Table~\ref{angular:tab:ten}. The first four rows
recover the source's theorem. The next six and the uniformity for
all $b>0$ are supplied here.

\begin{table}[ht]
\centering
\caption{The complete angular expansion above the scale $(2/13)^a$.
Equal rational bases have already been combined.}
\label{angular:tab:ten}
\begin{tabular}{@{}cl@{}}
\toprule
$q$& $c(q)$ in $T_{10}=\sum c(q)q^a$\\
\midrule
$2/3$ & $h_2/2$\\
$2/5$ & $-h_4/2$\\
$1/3$ & $h_2h_3$\\
$8/27$ & $-23h_2^3/48$\\
$2/7$ & $h_6/2$\\
$2/9$ & $-(3h_2h_5+h_8)/2$\\
$1/5$ & $-h_3h_4$\\
$2/11$ & $h_{10}/2$\\
$8/45$ & $39h_2^2h_4/16$\\
$1/6$ & $2h_2(h_3^2+h_7)$\\
\bottomrule
\end{tabular}
\end{table}

\begin{corollary}[A sharp signed remainder after ten terms]
\label{angular:cor:ten}
Uniformly for $b>0$,
\begin{equation}\label{angular:eq:eleven}
 \delta_{a,b}=T_{10}(a,b)-\frac{h_{12}}2(2/13)^a
       +O((4/27)^a).
\end{equation}
In particular,
\[
 \frac{\delta_{a,b}-T_{10}(a,b)}{(2/13)^a}
   =-\frac{h_{12}}2+O((26/27)^a).
\]
For all sufficiently large integer $a$, with a threshold independent
of $b>0$, the ten-term approximation strictly overestimates
$\delta_{a,b}$ and therefore gives a strict lower approximation to
$\theta_{a,b}$.
\end{corollary}

\begin{proof}
Apply Theorem~\ref{angular:thm:allorders} at the cutoff $4/27$.
The only additional nonzero base strictly between $4/27$ and
$2/13$, including $2/13$ itself, is $2/13$ with coefficient
$-h_{12}/2$.  This finite assertion can be checked by the multiset
bounds in the proof, or by the exact generator supplied with the
article.  The source four-term coefficient and all six new displayed
coefficients are also verified by substitution into the original
Fourier equation, independently of the Lagrange coefficient formula.
Finally $h_{12}\geq1$, so the uniform remainder gives the eventual
strict sign.
\end{proof}

For illustration, at $a=256$, $b=4$, direct Fourier root evaluation
gives
\[
 \delta_{256,4}\approx4.42523791578002118187450681399\times10^{-46}.
\]
The ten-term absolute error is approximately $4.24289006845\times10^{-209}$.
Its normalized error is $-0.541353663458730256$, approaching the
predicted limit
\[
 -H_{12}^{(4)}/2=-0.541076555073967515\ldots.
\]
The convergence of this last ratio is deliberately slow: its proved
next scale is $(26/27)^a$.  These decimal values are independent
high-precision diagnostics, not interval certificates or inputs to
the proof.  The supplied script bounds the omitted Fourier tail and
labels its floating-point status explicitly.

The next coefficient can also be made explicit. Extraction through the
cutoff $1/7$ gives
\begin{align*}
 \delta_{a,b}={}&T_{10}(a,b)-\frac{h_{12}}2(2/13)^a
 -\frac{27}{8}h_2^3h_3(4/27)^a+O((1/7)^a),
\end{align*}
uniformly for $b>0$. The nonzero contribution at base $4/27$ comes
from the multiset $(3,3,3,4)$; the other representations vanish by
parity. The accompanying audit extraction substitutes all twelve
terms back into the Fourier equation with exactly zero surviving
residual above the cutoff. This extra negative term explains the
direction of approach of the normalized numerical error.

\subsection{An exact and finite coefficient algorithm}

To compute the expansion above a prescribed rational cutoff $\eta$,
enumerate nondecreasing multisets of integers $n\geq3$ while
$\prod(2/n)>\eta$.  The total multiplicity and the possible indices
are bounded as in the theorem. Remove multisets containing an even
number of odd indices. Compute the coefficient in
\eqref{angular:eq:coefficient} with exact polynomial arithmetic and
combine equal rational bases. This algorithm terminates for every
cutoff, uses no integer-relation search, and returns the complete
coefficient set at that scale.

For the cutoff $2/13$, it examines 25 multisets and produces the ten
nonzero combined bases in Table~\ref{angular:tab:ten}. The verification
routine then substitutes the answer into \eqref{angular:eq:fourier}
in a separately implemented finite algebra of rational exponential
bases. Every surviving coefficient cancels exactly. Thus the artifact
checks the coefficient identities, while the theorem proves the
uniform analytic remainder and the applicability of those identities
to the actual zero.
